\section{Error-Free Computations}

Fix $\mathbb F=\mathbb F_{E, M}$ and $\rnd=\rnd_{E, M}$ as before.
For $a\in\mathbb F\setminus\{0\}$ let $e_a\in[e_{\min}, e_{\max}]\cap\mathbb Z$
denote its floating-point exponent, taking $e_a=e_{\min}$ for subnormals.
Set $e_0=e_{\min}$ when an exponent is needed for zero.
Write $\Delta_{\min}\coloneqq2^{e_{\min}-M}$ for the least positive
subnormal spacing. All exactness assertions below concern finite
results with the indicated range conditions.

\begin{lemma}[Integer Significands]
A real number $a$ is in $\mathbb F$ if and only if there exist
$m, e\in\mathbb Z$ with
\[
 a=m2^{e-M}, \qquad |m|<2^{M+1}, \qquad
 e_{\min}\le e\le e_{\max}.
\]
In particular, every $a\in\mathbb F$ is an integer multiple of
$\Delta_{\min}$.
\end{lemma}
\begin{proof}
For normal values take $e=e_a$ and let $m$ be the integer encoded
by the $(M+1)$ significand bits; for subnormals take $e=e_{\min}$
and $|m|<2^M$. Conversely, for any displayed pair $(m, e)$, a
significand with at most $M+1$ bits is either normal at a suitable
exponent or subnormal after shifting to $e_{\min}$. If a shift would
exceed $e_{\max}$, the bound $|m|<2^{M+1}$ prevents it.
\end{proof}

\begin{theorem}[Hauser: Exact Addition]
For $a, b\in\mathbb F$,
\[
 |a+b|\le2^{e_{\min}+1}\implies a\oplus b=a+b.
\]
\end{theorem}
\begin{proof}
Both $a$ and $b$ are integer multiples of $\Delta_{\min}$, hence
so is $a+b$. Every such multiple in
$[-2^{e_{\min}+1},2^{e_{\min}+1}]$ is representable. Therefore
$a+b\in\mathbb F$ and rounding does nothing.
\end{proof}

\begin{theorem}[Sterbenz: Exact Subtraction]
For $a, b\in\mathbb F$,
\[
 \frac a2\le b\le2a\implies a\ominus b=a-b.
\]
Equivalently, one may assume $b/2\le a\le2b$.
\end{theorem}
\begin{proof}
The hypothesis forces $a, b\ge0$. The zero case is immediate.
By exchanging $a, b$ and changing the sign if necessary, assume
$a\ge b$, so $b\le a\le2b$.
Write $a=m_a2^{e_a-M}$ and $b=m_b2^{e_b-M}$ by the lemma.
As $a\ge b$, we have $e_a\ge e_b$ (with the convention for
subnormals), and thus
\[
 a-b=(m_a2^{e_a-e_b}-m_b)2^{e_b-M}=m'2^{e_b-M},
 \qquad m'\in\mathbb Z.
\]
Since $0\le a-b\le b$, one has $0\le m'\le m_b<2^{M+1}$,
so $a-b\in\mathbb F$. Thus $a\ominus b=a-b$.
\end{proof}

\begin{proposition}[Exact Scaling by Powers of Two]
Suppose $a, b\in\mathbb F\setminus\{0\}$ have
$|a|=(\alpha_0.\alpha_1\dots\alpha_M)_2\,2^{e_a}$ and
$|b|=2^{e_b}$. If $e_a+e_b\in[e_{\min}, e_{\max}]$, then
\[
 a\otimes b=ab.
\]
Similarly, if $e_a-e_b\in[e_{\min}, e_{\max}]$, then
$a\odiv b=a/b$. \end{proposition}

\begin{proposition}[Exact Products with Trailing Zeros]
Let $a,b\in\mathbb F$ be nonzero normal values whose binary significands
have, respectively, at least $j$ and $j'$ trailing zero bits. If
\[
  j+j'\ge M+1,\qquad
  e_{\min}\le e_a+e_b\le e_{\max}-1,
\]
then $ab\in\mathbb F$ and $a\otimes b=ab$.
\end{proposition}
\begin{proof}
Removing the trailing zeros leaves integer significands of at most
$M+1-j$ and $M+1-j'$ bits. Their product has at most
$2(M+1)-j-j'\le M+1$ bits. The exact product therefore fits the
available significand precision; its normalized exponent is either
$e_a+e_b$ or $e_a+e_b+1$, both representable by hypothesis.
\end{proof}

\begin{remark}
Sterbenz's theorem does not preclude catastrophic cancellation:
subtraction of nearby \emph{represented} floats may be exact, while
an error inherited from earlier rounded operations becomes large
relative to the small difference.
\end{remark}

\begin{theorem}[Dekker's Theorem (Fast2Sum)]
For $a, b\in\mathbb F$ with $|a+b|\le F_{\max}$, put
\[
 s=a\oplus b, \qquad r=(a+b)-s.
\]
Then $r\in\mathbb F$. If $e_a\ge e_b$, the algorithm
\[
 \begin{aligned}
 s&=a\oplus b, \\
 b'&=s\ominus a, \\
 z&=b\ominus b'
 \end{aligned}
 \qquad\text{(Fast2Sum)}
\]
returns $(s, z)=(a\oplus b, r)$; in particular, $a+b=s+z$
\emph{exactly}.
\end{theorem}
\begin{proof}
By symmetry assume $e_a\ge e_b$. We first show $s-a\in\mathbb F$.
If $a,b$ have the same sign and $|a|\ge|b|$, then $s$ has that sign
and $|a|\le|s|\le2|a|$; Sterbenz's theorem applies to $s-a$.
If their signs differ and $|b|\ge|a|/2$, their magnitudes differ
by at most a factor of two, so Sterbenz makes $a+b$ exact.
If their signs differ and $|b|<|a|/2$, either the sum is exact
in the subnormal range by Hauser's theorem, or $s$ and $a$ are
within a factor of two and Sterbenz applies again.
In the remaining case, $a,b$ have the same sign with
$|b|>|a|$, which forces $e_a=e_b$. Both belong to the same
binade of spacing $h=2^{e_a-M}$; $s-a$ is a multiple of $h$
with magnitude at most $2^{e_a+1}$, hence is representable.
Thus $b'=s\ominus a=s-a$.

Moreover, $a$, $b$, and $s$ lie on the grid $2^{e_b-M}\mathbb Z$;
if the exact sum is smaller than $2^{e_b}$, it is itself
representable, and otherwise rounding uses a no-finer grid.
Consequently, $r=b-b'$ is an integer multiple of $2^{e_b-M}$.
By optimality of nearest rounding, comparison with the
representable candidate $a$ yields
\[
 |r|=|(a+b)-s|\le|(a+b)-a|=|b|.
\]
The integer-significand lemma now implies $r\in\mathbb F$.
Therefore $b\ominus b'=r$, giving $a+b=s+r$ exactly.
\end{proof}

\begin{theorem}[Boldo's Theorem (Product Residual)]
For $a, b\in\mathbb F$ with $|ab|\le F_{\max}$, let
$s=a\otimes b$ and $r=ab-s$. If
\[
 e_a+e_b\ge e_{\min}+M,
\]
then $r\in\mathbb F$ and
\[
 r=\fma(a, b,-s).
\]
\end{theorem}
\begin{proof}
In the integer-significand representation,
$ab=m_am_b2^{e_a+e_b-2M}$. The rounded product $s$ lies on the same
binary grid, so
$ab-s=k2^{e_a+e_b-2M}$ for some integer $k$. The exact product
has exponent $e_a+e_b$ or $e_a+e_b+1$; nearest rounding therefore
gives $|k|\le2^M$. By the hypothesis,
$e_a+e_b-M\ge e_{\min}$, so the residual has a representable
significand and exponent. Thus $r\in\mathbb F$, and the
single-rounding property of FMA gives
$\fma(a,b,-s)=\rnd(r)=r$.
\end{proof}


